%! Equivalent Representations of Trigonometric Functions — Unit 7, Lesson 7.5 — Lesson Plan
\documentclass[10pt]{article}
\usepackage{precalculus-article}
\usepackage{precalculus-boxes}
\usepackage{graphicx}
\graphicspath{{images/}}
\newcommand{\TallMath}[1]{$\displaystyle #1\rule[-1.4em]{0pt}{3.2em}$}
\newcommand{\CourseName}{Precalculus}
\newcommand{\MeetingLength}{60 minutes}
\newcommand{\UnitNumberName}{Unit 7: Analytic Trigonometry and Polar Coordinates \quad}
\newcommand{\LessonNumberName}{Lesson 7.5: Equivalent Representations of Trigonometric Functions}

\begin{document}

%----------------------------------------------------------------------
% Title Block
%----------------------------------------------------------------------
\begin{center}
  {\Large\bfseries \CourseName} \\[4pt]
  {\small\bfseries \UnitNumberName \LessonNumberName \quad (3--4 instructional periods, \MeetingLength\ each)}
\end{center}

%----------------------------------------------------------------------
% Primary Objective
%----------------------------------------------------------------------
\begin{tcolorbox}[colback=sky, colframe=navy, boxrule=0.9pt, arc=2mm,
    left=2mm, right=2mm, top=1.5mm, bottom=1.5mm]
  \textbf{Primary Objective:} Students will use the Pythagorean identity
  $\sin^2\theta + \cos^2\theta = 1$ and its derived forms, together with the
  sum, difference, and double-angle identities for sine and cosine, to rewrite
  trigonometric expressions in equivalent forms and to solve trigonometric
  equations analytically.
  \textit{(Big Idea: Functions)}
\end{tcolorbox}

\vspace{0.08in}

%----------------------------------------------------------------------
% Priority Ideas & Skills
%----------------------------------------------------------------------
\begin{skillbox}[Priority Ideas \& Skills]{goldbox}
\begin{tabularx}{\linewidth}{@{} p{0.46\linewidth} X @{}}
\textbf{AP Skills} & \textbf{Key Understandings (EKs)} \\
\midrule
\textbf{1.A} \textit{Solve equations/inequalities analytically} ---
  Use algebraic manipulation of trig identities to transform equations
  into solvable forms (e.g., convert $2\sin^2\theta + \cos\theta - 1 = 0$
  to a quadratic in $\cos\theta$).
&
\textbf{EK 3.12.A.1}: The Pythagorean Theorem applied to the unit circle
  gives $\sin^2\theta + \cos^2\theta = 1$ (the Pythagorean identity). \\[4pt]

\textbf{1.B} \textit{Express in equivalent forms} ---
  Rewrite trig expressions using the Pythagorean identity and its variants
  ($1 + \tan^2\theta = \sec^2\theta$, $\cot^2\theta + 1 = \csc^2\theta$),
  and using sum/difference/double-angle identities.
&
\textbf{EK 3.12.A.2}: Algebraic manipulation of the Pythagorean identity
  produces $\tan^2\theta = \sec^2\theta - 1$, $\cot^2\theta = \csc^2\theta - 1$,
  and relationships such as $\arcsin x = \arccos\!\sqrt{1 - x^2}$ (with
  appropriate domain restrictions). \\[4pt]

\textbf{3.B} \textit{Apply numerical results} ---
  Use identities to compute exact values of trig functions at non-standard
  angles (e.g., $\sin 75^\circ = \sin(45^\circ + 30^\circ)$) and interpret
  results in context.
&
\textbf{EK 3.12.B.1--B.4}: The sum identities
  $\sin(\alpha+\beta) = \sin\alpha\cos\beta + \cos\alpha\sin\beta$ and
  $\cos(\alpha+\beta) = \cos\alpha\cos\beta - \sin\alpha\sin\beta$ extend
  to difference and double-angle forms, and can be used to verify additional
  identities. \\[4pt]

&
\textbf{EK 3.12.C.1--C.2}: Choosing the right equivalent form makes
  information more accessible; equivalent trig forms are especially useful
  for solving equations and inequalities. \\
\end{tabularx}
\end{skillbox}

\vspace{0.08in}

%----------------------------------------------------------------------
% Vocabulary
%----------------------------------------------------------------------
\begin{skillbox}[{Vocabulary, Concepts, \& Theorems}]{greenbox}
\begin{tabularx}{\linewidth}{@{} p{0.28\linewidth} X @{}}
\textbf{Term} & \textbf{Definition / Note} \\
\midrule
Pythagorean identity &
  $\sin^2\theta + \cos^2\theta = 1$ --- derived from the unit-circle
  coordinates $(\cos\theta, \sin\theta)$ and the Pythagorean Theorem. \\[4pt]
sum identity (sine) &
  $\sin(\alpha+\beta) = \sin\alpha\cos\beta + \cos\alpha\sin\beta$. \\[4pt]
sum identity (cosine) &
  $\cos(\alpha+\beta) = \cos\alpha\cos\beta - \sin\alpha\sin\beta$. \\[4pt]
difference identity &
  Replace $\beta$ with $-\beta$ in the sum identity and apply even/odd
  function properties to obtain formulas for $\sin(\alpha-\beta)$ and
  $\cos(\alpha-\beta)$. \\[4pt]
double-angle identity &
  Setting $\alpha = \beta$ in the sum identity gives
  $\sin 2\theta = 2\sin\theta\cos\theta$ and
  $\cos 2\theta = \cos^2\theta - \sin^2\theta$. \\[4pt]
verify an identity &
  Show that two trig expressions are equal for all values in their domains
  by transforming one side algebraically (work from one side only). \\
\end{tabularx}
\end{skillbox}

\vspace{0.08in}

%----------------------------------------------------------------------
% Activate Prior Knowledge
%----------------------------------------------------------------------
\begin{skillbox}[Activate Prior Knowledge \& Spiral Review]{sky}
\begin{tabularx}{\linewidth}{@{}X c@{}}
\begin{minipage}[t]{0.96\linewidth}\vspace{0pt}
\textbf{Spiral topics activated during Warm-Up (first 8--10 minutes):}
\begin{itemize}[leftmargin=1.4em, itemsep=2pt]
  \item \textbf{Unit-circle exact values}: $\sin\theta$ and $\cos\theta$ at
    $\pi/6$, $\pi/4$, $\pi/3$, $\pi/2$ --- students verify $\sin^2\theta + \cos^2\theta = 1$
    numerically before seeing it as a general identity.
  \item \textbf{Factoring quadratics and the quadratic formula}: The equation
    $2\cos^2\theta + \cos\theta - 1 = 0$ is quadratic in $\cos\theta$; fluency
    with factoring $(2u - 1)(u + 1) = 0$ is essential for Lesson 7.5.C.
  \item \textbf{Algebraic expansion}: $(a+b)^2 = a^2 + 2ab + b^2$ primes
    students to recognize the structure of the Pythagorean identity.
  \item \textbf{Ratio definitions}: $\tan\theta = \sin\theta / \cos\theta$,
    $\sec\theta = 1/\cos\theta$ from Lessons 3.9--3.11 --- students need these
    to derive the second Pythagorean form.
\end{itemize}
\end{minipage}
&
\end{tabularx}
\end{skillbox}

\vspace{0.08in}

%----------------------------------------------------------------------
% Hook
%----------------------------------------------------------------------
\begin{skillbox}[Hook --- Entry Event]{sky}
\textbf{Scenario:} Write the following on the board: $2\sin^2\theta + \cos\theta = 1$.
Ask students to try solving for $\theta$ directly. After a minute of struggle,
ask: \textit{``Both sine and cosine appear. How could you rewrite this with
only \emph{one} trig function?''}

\textbf{Discussion prompts:}
\begin{enumerate}[leftmargin=1.4em, itemsep=3pt]
  \item \textit{``We know $\sin^2\theta + \cos^2\theta = 1$ from the unit
    circle. Can you rearrange it to express $\sin^2\theta$ in terms of
    $\cos\theta$?''}
  \item \textit{``If you substitute that expression, what kind of equation
    do you get? Can you solve it?''}
  \item \textit{``What other relationships might we derive from that single
    identity?''}
\end{enumerate}

\textbf{Key tension to surface:} Trig equations involving more than one
function often cannot be solved directly. An identity --- an algebraic rule
that is always true --- lets us rewrite every term using a single function,
converting the problem into a solvable algebraic equation.

\textbf{Transition:} ``Today we'll build a toolkit of trig identities ---
starting from the one identity that comes straight from the unit circle ---
and use them to crack equations that would otherwise be impossible to solve.''
\end{skillbox}

\vspace{0.08in}

%----------------------------------------------------------------------
% Lesson
%----------------------------------------------------------------------
\begin{skillbox}[Lesson --- Instruction Sequence]{sky}
\begin{multicols}{2}

\textbf{Step 1 --- Pythagorean Identity (LO 3.12.A)} (12 min)

Draw the unit circle; mark a point $(\cos\theta, \sin\theta)$.
The distance from origin to this point is 1 (radius), so by the
Pythagorean Theorem: $\cos^2\theta + \sin^2\theta = 1$.
This is \emph{always} true, for every $\theta$ --- it is an \emph{identity},
not just an equation.

Derive two more forms: divide through by $\cos^2\theta$ to get
$1 + \tan^2\theta = \sec^2\theta$; divide through by $\sin^2\theta$ to get
$\cot^2\theta + 1 = \csc^2\theta$. Students complete fill-in derivation in notes.

\columnbreak

\textbf{Step 2 --- Sum and Difference Identities (LO 3.12.B)} (14 min)

State the formulas (do not prove algebraically at this level):
$\sin(\alpha+\beta) = \sin\alpha\cos\beta + \cos\alpha\sin\beta$
$\cos(\alpha+\beta) = \cos\alpha\cos\beta - \sin\alpha\sin\beta$

Worked example: find $\sin 75^\circ$ exactly.
$\sin(45^\circ+30^\circ)=\sin 45^\circ\cos 30^\circ + \cos 45^\circ\sin 30^\circ$
$=\dfrac{\sqrt{2}}{2}\cdot\dfrac{\sqrt{3}}{2}+\dfrac{\sqrt{2}}{2}\cdot\dfrac{1}{2}
=\dfrac{\sqrt{6}+\sqrt{2}}{4}$

Then derive difference identities by replacing $\beta$ with $-\beta$ and
using $\sin(-\beta)=-\sin\beta$, $\cos(-\beta)=\cos\beta$.

\end{multicols}

\begin{multicols}{2}

\textbf{Step 3 --- Double-Angle Identities (LO 3.12.B)} (10 min)

Set $\alpha = \beta = \theta$ in the sum identities:
$\sin 2\theta = 2\sin\theta\cos\theta$
$\cos 2\theta = \cos^2\theta - \sin^2\theta$

Alternate cosine forms (by Pythagorean substitution):
$\cos 2\theta = 2\cos^2\theta - 1 = 1 - 2\sin^2\theta$

Useful rearrangements for solving: $\cos^2\theta = \frac{1+\cos 2\theta}{2}$,
$\sin^2\theta = \frac{1 - \cos 2\theta}{2}$.

\columnbreak

\textbf{Step 4 --- Solving Equations with Identities (LO 3.12.C)} (12 min)

\textbf{Example A:} $2\sin^2\theta + \cos\theta - 1 = 0$.
Substitute $\sin^2\theta = 1 - \cos^2\theta$:
$2(1-\cos^2\theta)+\cos\theta-1=0 \Rightarrow -2\cos^2\theta + \cos\theta + 1 = 0$
$\Rightarrow 2\cos^2\theta - \cos\theta - 1 = 0 \Rightarrow (2\cos\theta+1)(\cos\theta-1)=0$
Solutions in $[0,2\pi)$: $\theta = 0,\; 2\pi/3,\; 4\pi/3$.

\textbf{Example B:} Verify $\tan\theta\cos\theta = \sin\theta$.
Left side: $\dfrac{\sin\theta}{\cos\theta}\cdot\cos\theta = \sin\theta$ = right side.

\end{multicols}

\smallskip
\textbf{Pacing note:} Steps 1--2 in Period 1; Steps 3--4 in Period 2;
activity and exit ticket in Period 3 (Period 4 if re-teach needed).
\end{skillbox}

\vspace{0.08in}

%----------------------------------------------------------------------
% Active Monitoring
%----------------------------------------------------------------------
\begin{skillbox}[Active Monitoring]{redbox}
\begin{itemize}[leftmargin=1.4em, itemsep=3pt]
  \item \textbf{Squaring both sides illegally}: Students may try to isolate
    $\sin\theta$ and square, introducing extraneous solutions. Prompt:
    \textit{``Did you check every solution back in the original equation?''}
  \item \textbf{Wrong substitution direction}: Students substitute
    $\cos^2\theta = 1 + \sin^2\theta$ (incorrect sign). Reinforce:
    $\cos^2\theta = 1 - \sin^2\theta$ (the $\sin^2\theta$ term moves with a minus sign).
  \item \textbf{Confusing sum identity sign}: For cosine, the sum identity has a
    \emph{minus}: $\cos(\alpha+\beta) = \cos\alpha\cos\beta - \sin\alpha\sin\beta$.
    Students often write $+$. Have them compare with the sine sum identity.
  \item \textbf{``Verifying'' by plugging in numbers}: Remind students that a
    numerical check does not constitute a proof of an identity; they must work
    algebraically from one side to the other.
  \item \textbf{Forgetting all solutions in $[0, 2\pi)$}:
    After solving $\cos\theta = -1/2$, students may give only $2\pi/3$ and miss $4\pi/3$.
    Prompt: \textit{``In which quadrants is cosine negative? There should be two solutions.''}
\end{itemize}
\end{skillbox}

\vspace{0.08in}

%----------------------------------------------------------------------
% Group Work & Differentiation
%----------------------------------------------------------------------
\begin{skillbox}[Group Work \& Differentiation]{redbox}
Students work in groups of 3--4 on the tiered activity (teacher-assigned
or student self-selected with guidance).

\begin{itemize}[leftmargin=1.4em, itemsep=4pt]
  \item \textbf{Tier R (Reinforce)}: Given the equation $2\sin^2x + \cos x - 1 = 0$
    and the substitution $\sin^2x = 1 - \cos^2x$ already provided, students complete
    each scaffolded algebraic step, factor the resulting quadratic, and solve.
    Answer-checking table provided.
  \item \textbf{Tier A (Apply)}: Three trig equations, each requiring a different
    identity strategy (Pythagorean substitution, double-angle, purely factored).
    Students solve for all solutions in $[0, 2\pi)$ without scaffolded steps.
  \item \textbf{Tier E (Extend)}: Students verify two trig identities (one side only),
    then use a sum identity to find an exact value of $\sin(5\pi/12)$.
    Optional: prove a novel identity given a hint about which identity to apply first.
\end{itemize}
\end{skillbox}

\vspace{0.08in}

%----------------------------------------------------------------------
% Individual Work & Assessment
%----------------------------------------------------------------------
\begin{skillbox}[Individual Work \& Assessment]{redbox}
Exit ticket administered independently (final 5--7 minutes of Period 3).

\begin{enumerate}[leftmargin=1.4em, itemsep=4pt]
  \item Express $\sin^2\theta$ in terms of $\cos 2\theta$ (double-angle form).
  \item Use the difference identity to find the exact value of $\cos(\pi/12)$.
  \item Solve $\cos^2\theta - \sin^2\theta = 0$ for $\theta \in [0, 2\pi)$,
    using the double-angle identity to simplify first.
\end{enumerate}

\textbf{Data use:} Sort exit tickets by error pattern --- sign errors in sum/difference
identities, incorrect substitution directions, missed solutions due to quadrant errors.
Use distribution to plan a brief targeted opener for Lesson 3.13.
\end{skillbox}

\vspace{0.08in}

%----------------------------------------------------------------------
% Reinforcement & Extension
%----------------------------------------------------------------------
\begin{skillbox}[Reinforcement \& Extension]{goldbox}
\textbf{Homework (Lesson 7.5):} 5--6 problems covering: simplification using
the Pythagorean identity; exact-value computation using sum/difference identities;
one verify-the-identity problem (work from one side only); solving a trig equation
requiring identity substitution; 2 AP-style multiple-choice items.

\textbf{Preview of Lesson 3.13 --- Trigonometric Equations and Inequalities:}
Ask students: \textit{``You just solved equations like $2\sin^2\theta + \cos\theta - 1 = 0$.
What if the equation were $\cos 2\theta \geq \tfrac{1}{2}$? In Lesson 3.13 we'll extend
from equations to inequalities, and from specific intervals to general solutions in terms
of $2\pi k$.''}
\end{skillbox}

\end{document}
